Choose an affine open neighborhood $V$ of $P$ and put $U=f^{-1}(V)$. Restriction to the completion along $Y$ gives $\Gamma(U,\mathcal O_U)\to k$ by \exref{II.9.1}. This map is injective: an element in its kernel has germ in $\bigcap_{r\ge1}\mathcal I_y^r=0$ at any $y\in Y$, by the Krull intersection theorem (Atiyah--Macdonald, \emph{Introduction to Commutative Algebra}, Corollary 10.19), and therefore vanishes on the integral scheme $U$. Consequently every regular function on $V$ pulls back to its value at $P$, so $f|_U$ is constant. The closed subset $f^{-1}(P)$ contains the dense open subset $U$, hence equals $X$.
